\title{Rust 异步运行时设计}
\author{杨子凡}
\date{Oct 02, 2026}
\maketitle
Rust 异步编程的流行并非偶然。在高并发场景下，线程模型的上下文切换开销与 C10k 瓶颈使开发者不得不转向事件驱动架构。本文将带你手写一个极简的异步运行时，理解 Executor、Reactor 与调度策略的核心思想，并对比 tokio、async-std 与 smol 三种成熟方案。\par
\chapter{背景：为什么需要异步运行时}
当我们使用传统线程来处理并发请求时，每个连接往往需要一个线程来维持。线程的创建、调度与销毁都涉及昂贵的内核态上下文切换。面对数万并发连接时，系统可能耗尽内存或因频繁切换而导致吞吐量下降。Rust 虽然提供了 \texttt{std::thread}、\texttt{std::sync::mpsc} 与 \texttt{Mutex} 等原语，但在高并发下这些同步工具容易成为瓶颈。异步的本质是用「状态机 + 事件循环」来替代阻塞等待，让单线程或少量线程就能驱动大量并发任务。\par
\chapter{异步运行时核心抽象}
\section{Future 与 Poll}
在 Rust 中，所有异步操作都表现为实现了 \texttt{Future} trait 的类型。\texttt{Future} 的核心方法是 \texttt{poll}，它接收一个 \texttt{Context} 并返回 \texttt{Poll<T>}。当 \texttt{poll} 返回 \texttt{Poll::Ready(val)} 时，表明异步操作已完成；返回 \texttt{Poll::Pending} 则表示尚未就绪，需要等待外部事件唤醒。理解这一点对后续实现 Executor 至关重要。\par
\begin{lstlisting}[language=rust]
use std::future::Future;
use std::pin::Pin;
use std::task::{Context, Poll};

pub struct MyFuture { /* 内部状态 */ }

impl Future for MyFuture {
    type Output = i32;

    fn poll(self: Pin<&mut Self>, cx: &mut Context<'_>) -> Poll<Self::Output> {
        // 1. 检查内部状态是否就绪
        if self.ready {
            // 2. 返回就绪值
            Poll::Ready(42)
        } else {
            // 3. 注册唤醒器，等待外部事件
            cx.waker().wake_by_ref();
            Poll::Pending
        }
    }
}
\end{lstlisting}
这段代码演示了 \texttt{poll} 的基本流程：检查状态、返回结果或挂起。注意，\texttt{wake\_{}by\_{}ref} 会通知运行时该任务可再次被调度。\par
\section{Waker 与任务唤醒}
\texttt{Waker} 是连接 Future 与 Executor 的桥梁。当 Reactor 检测到 I/O 事件时，它会调用与该事件关联的 \texttt{Waker} 来唤醒对应任务。\texttt{Waker} 的实现通常包含一个指向任务队列的引用或索引，以便 Executor 能迅速找到待执行任务。\par
\section{Executor 与 Reactor 的职责划分}
Executor 负责管理任务队列与调度策略，而 Reactor 则专注于监听 I/O 事件并通过多路复用机制（如 epoll、kqueue 或 IOCP）高效分发。两者通过 \texttt{Waker} 进行协作：Reactor 产生事件，Executor 消费事件并驱动 Future 状态机。\par
\chapter{手写一个最小运行时}
\section{项目结构与依赖}
我们先创建一个仅依赖 \texttt{std} 与 \texttt{libc} 的项目，命名为 \texttt{minirt}。目录结构如下：\par
\begin{lstlisting}
minirt/
├── Cargo.toml
├── src/
│   ├── main.rs
│   ├── executor.rs
│   └── reactor.rs
\end{lstlisting}
\texttt{Cargo.toml} 中仅声明 \texttt{libc} 以便调用系统调用。\par
\section{实现极简的 block\_{}on}
\texttt{block\_{}on} 是运行时的入口，它会驱动一个 Future 直到完成。\par
\begin{lstlisting}[language=rust]
use std::future::Future;
use std::pin::Pin;
use std::task::{Context, Poll, RawWaker, RawWakerVTable, Waker};
use std::sync::Arc;

pub fn block_on<F: Future>(mut fut: F) -> F::Output {
    // 1. 创建一个简单的 Waker，它什么也不做
    let waker = unsafe { Waker::from_raw(dummy_raw_waker()) };
    let mut cx = Context::from_waker(&waker);
    // 2. 将 Future 固定在栈上
    let mut fut = unsafe { Pin::new_unchecked(&mut fut) };
    // 3. 循环调用 poll，直到 Ready
    loop {
        match fut.as_mut().poll(&mut cx) {
            Poll::Ready(val) => return val,
            Poll::Pending => {
                // 4. 在单线程示例中，直接让出 CPU
                std::thread::yield_now();
            }
        }
    }
}

// 省略 dummy_raw_waker 的实现
\end{lstlisting}
这段代码展示了 \texttt{block\_{}on} 的核心逻辑：不断调用 \texttt{poll}，直到 Future 返回就绪值。在真实场景中，\texttt{yield\_{}now} 会被更高效的事件等待替代。\par
\section{任务队列}
Executor 需要维护一个任务队列。单线程版本可使用 \texttt{VecDeque}，跨线程版本则需使用 \texttt{crossbeam::deque} 或 \texttt{std::sync::mpsc}。下面是一个基于 \texttt{VecDeque} 的简单实现：\par
\begin{lstlisting}[language=rust]
use std::collections::VecDeque;
use std::future::Future;
use std::pin::Pin;
use std::task::{Context, Poll};

pub struct Executor {
    tasks: VecDeque<Pin<Box<dyn Future<Output = ()>>>>,
}

impl Executor {
    pub fn new() -> Self {
        Executor { tasks: VecDeque::new() }
    }

    pub fn spawn<F: Future<Output = ()> + 'static>(&mut self, fut: F) {
        self.tasks.push_back(Box::pin(fut));
    }

    pub fn run(&mut self) {
        while let Some(mut task) = self.tasks.pop_front() {
            // 创建一个 dummy waker，实际使用时需替换
            let waker = unsafe { Waker::from_raw(dummy_raw_waker()) };
            let mut cx = Context::from_waker(&waker);
            if task.as_mut().poll(&mut cx) == Poll::Pending {
                self.tasks.push_back(task);
            }
        }
    }
}
\end{lstlisting}
\texttt{spawn} 将 Future 装箱后推入队列，\texttt{run} 则依次取出并驱动，直到所有任务完成或再次挂起。\par
\section{Reactor：用 epoll 监听事件}
Reactor 的职责是监听文件描述符上的 I/O 事件。以 Linux 为例，我们使用 \texttt{epoll}。\par
\begin{lstlisting}[language=rust]
use libc::{self, epoll_event, EPOLLIN};
use std::os::unix::io::RawFd;

pub struct Reactor {
    epoll_fd: RawFd,
}

impl Reactor {
    pub fn new() -> std::io::Result<Self> {
        // 1. 创建 epoll 实例
        let epoll_fd = unsafe { libc::epoll_create1(0) };
        if epoll_fd < 0 {
            return Err(std::io::Error::last_os_error());
        }
        Ok(Reactor { epoll_fd })
    }

    pub fn register(&self, fd: RawFd, waker: Waker) -> std::io::Result<()> {
        // 2. 将 fd 及其关联的 waker 注册到 epoll
        let mut ev = epoll_event {
            events: EPOLLIN as u32,
            u64: Box::into_raw(Box::new(waker)) as u64,
        };
        let ret = unsafe {
            libc::epoll_ctl(self.epoll_fd, libc::EPOLL_CTL_ADD, fd, &mut ev)
        };
        if ret < 0 {
            Err(std::io::Error::last_os_error())
        } else {
            Ok(())
        }
    }

    pub fn poll(&self, timeout_ms: i32) -> std::io::Result<Vec<Waker>> {
        // 3. 等待事件，最多返回 128 个
        let mut events = [epoll_event { events: 0, u64: 0 }; 128];
        let n = unsafe {
            libc::epoll_wait(self.epoll_fd, events.as_mut_ptr(), 128, timeout_ms)
        };
        if n < 0 {
            return Err(std::io::Error::last_os_error());
        }
        // 4. 提取 waker 并唤醒
        let mut wakers = Vec::new();
        for ev in &events[..n as usize] {
            let waker = unsafe { Box::from_raw(ev.u64 as *mut Waker) };
            waker.wake();
            wakers.push(*waker);
        }
        Ok(wakers)
    }
}
\end{lstlisting}
\texttt{register} 将文件描述符与 \texttt{Waker} 绑定，\texttt{poll} 则阻塞等待事件并唤醒对应任务。\par
\section{把 Future 状态机与 epoll 事件关联}
当一个异步 TCP 连接需要等待数据时，Future 会把自身的 \texttt{Waker} 注册到 Reactor。数据到达后，Reactor 通过 epoll 事件唤醒 Future，Future 再次被 Executor 调度。\par
\section{一次完整的 echo server 请求生命周期}
\begin{itemize}
\item 服务器启动时创建监听 socket 并注册到 Reactor。
\item 客户端连接到达，Reactor 产生 \texttt{EPOLLIN} 事件，唤醒 \texttt{accept} Future。
\item \texttt{accept} Future 返回新连接，Executor 将 echo Future 加入队列。
\item echo Future 尝试 \texttt{read}，若无数据则把 \texttt{Waker} 注册到 Reactor 并返回 \texttt{Pending}。
\item 数据到达，Reactor 唤醒 echo Future，Future 继续执行并写回响应。
\end{itemize}
\chapter{调度策略与性能权衡}
\section{Work-stealing 与 Global queue}
Work-stealing 调度器为每个线程维护一个本地队列，空闲线程可从其他线程「偷取」任务，降低全局锁竞争。Global queue 则所有线程共享一个队列，简单但易成为瓶颈。\par
\section{优先级与饥饿}
为避免低优先级任务长期得不到执行，调度器通常引入优先级或 aging 机制。tokio 通过注入「LIFO 插队」与「公平轮转」来平衡响应延迟与吞吐量。\par
\section{批量 poll 与 spin-loop 阈值}
连续多次 \texttt{poll} 可能导致 CPU 空转。运行时会设置一个「spin-loop」阈值，超过后让出 CPU 或进入休眠，等待 Reactor 事件。\par
\section{NUMA 与 CPU 亲和性}
在 NUMA 架构下，跨节点访问内存代价高昂。tokio 提供了 \texttt{tokio::runtime::Builder::on\_{}thread\_{}park} 等钩子，允许用户设置 CPU 亲和性，减少跨节点通信。\par
\chapter{成熟运行时对比}
tokio 采用 work-stealing 调度与 mio 后端，拥有最成熟的生态与零拷贝支持。async-std 同样使用 work-stealing，但更强调与标准库一致的 API。smol 则定位为轻量级方案，默认单线程，适合嵌入式或教学场景。我们手写的 minirt 仅实现 FIFO 与 epoll，缺乏定时器与零拷贝，生态成熟度最低，但足以展示核心原理。\par
\chapter{工程化考量}
\section{取消与超时}
\texttt{select!} 宏可同时驱动多个 Future，任意一个完成即取消其余。\texttt{timeout} 则在指定时间内未完成即返回错误，二者都依赖 \texttt{Waker} 的正确注册与注销。\par
\section{背压与流量控制}
异步 channel 的容量直接影响背压。若发送方速度远快于接收方，队列可能耗尽内存。合理设置容量并在满时阻塞发送方，可避免系统雪崩。\par
\section{内存模型与 Pin}
\texttt{Pin} 保证 Future 在 \texttt{poll} 期间不会被移动，从而安全使用自引用结构。开发者需谨慎使用 \texttt{Box::pin} 与 \texttt{Pin::new\_{}unchecked}，避免悬垂指针。\par
\section{调试}
\texttt{tracing} crate 可记录任务生命周期，tokio-console 提供实时任务视图，火焰图则帮助定位 CPU 热点。\par
\chapter{案例研究：把最小运行时扩展为支持多线程}
\section{引入全局/局部队列}
多线程版本为每个线程分配一个本地 \texttt{Worker}，并共享一个全局 \texttt{Injector}。当本地队列为空时，线程先尝试从全局队列窃取任务，再尝试从其他线程窃取。\par
\section{无锁并发数据结构}
\texttt{crossbeam::deque} 提供无锁双端队列，支持 \texttt{steal} 与 \texttt{push} 的并发操作，避免了全局锁。\par
\section{线程池初始化与优雅退出}
在 \texttt{main} 中创建固定数量的线程，每个线程运行 \texttt{Worker::run}。当收到关闭信号时，线程会完成当前任务后退出，确保资源正确释放。\par
\chapter{未来与生态}
\section{io\_{}uring 与高性能内核旁路}
io\_{}uring 允许用户态与内核通过共享内存提交与完成 I/O 请求，减少系统调用开销。Rust 生态已有 \texttt{io\_{}uring} crate，可作为 Reactor 的高性能后端。\par
\section{异步与 async fn 语法稳定化}
\texttt{async fn} 与 \texttt{await} 已成为稳定语法，Future 状态机由编译器自动生成，极大降低了异步编程的心智负担。\par
\section{跨语言运行时互操作}
C\#{} 的 \texttt{Task}、Python 的 \texttt{asyncio} 与 Rust 的 \texttt{Future} 可通过 FFI 或消息队列实现互操作，构建混合语言的高并发服务。\par
通过手写 minirt，我们理解了 Executor、Reactor 与调度策略的协作方式。下一步可尝试实现带优先级的调度、支持超时的 \texttt{timeout} 以及接入 io\_{}uring，以进一步提升性能。推荐阅读 The tokio-rs book、Programming Rust 第 2 版中的「Rust Concurrency」章节，以及 epoll 与 io\_{}uring 的官方文档。\par
